\section{Conclusion}{\label{sec:conclusion}}

Systems code must adhere to stringent requirements on data
representation to achieve efficient, predictable performance and avoid
costly mediation at abstraction boundaries. In many cases, these
requirements result in code that is error prone and tedious to write,
ugly to read, and very difficult to verify.

We are not without hope, however, as we have demonstrated in this
paper. By using \Dargent, we can avoid the need for having the
\emph{glue} code (be it manually written or synthesised) that marshals
data from one format into another, and eliminate error-prone
bit-twiddling operations for manipulating specific bits in device
registers. Instead, we enable programmers to provide declarative
specifications of how their algebraic datatypes are laid out. Given
these specifications, our \emph{certifying} compiler generates
corresponding C code that operates on these data types directly,
 along with proofs that the generated code is
functionally correct. We have shown the applicability of \dargent on a
number of examples showcasing its support for low-level systems
features including the formal verification of a timer device driver.

